\documentclass[12pt]{article}
\usepackage[margin=0.6in,top=0.85in,bottom=0.55in]{geometry}
\usepackage{lmodern}
\usepackage{microtype}
\usepackage{fancyhdr}
\usepackage{titlesec}
\usepackage{enumitem}
\usepackage{lastpage}
\usepackage[hidelinks]{hyperref}

\setlength{\parindent}{0pt}
\setlength{\parskip}{0.25em}
\setlength{\headheight}{26pt}
\titleformat{\section}{\large\bfseries\scshape}{}{0pt}{}
\pagestyle{fancy}
\fancyhf{}
\fancyhead[C]{\footnotesize Student Name: \hspace{2.3in} \quad Assignment ID: U1L5LE \quad Page \thepage\ of \pageref{LastPage}}
\renewcommand{\headrulewidth}{0.5pt}

\newcommand{\answerbox}[2]{%
\vspace{0.12em}
\noindent\textbf{#1}\par
\vspace{0.04em}
\noindent\fbox{%
\begin{minipage}[t][#2]{\dimexpr\textwidth-2\fboxsep-2\fboxrule}
\rule{\textwidth}{0pt}
\vspace{#2}
\end{minipage}}
\par\vspace{0.22em}}

\begin{document}

\begin{center}
{\LARGE\bfseries Week 5 Lit Example Worksheet}\par\vspace{0.05em}
{\large\scshape Julius Caesar: Brutus, Antony, and Formal Shape}\par\vspace{0.12em}
\rule{0.78\textwidth}{0.5pt}
\end{center}

\textbf{Purpose:} Apply Week 5's Whately method to Julius Caesar. Identify mood and formal validity before judging factual strength.

\textbf{Directions:} Use the Lit Example Reader. Your answers should distinguish form test from evidence test.

\section*{Part 1: Brutus's Form}

\textbf{Part 1 Q1:} Reconstruct Brutus's argument in three claims: major premise, minor premise, and conclusion.

\answerbox{Part 1 Q1 ANSWER BOX}{2.45in}

\textbf{Part 1 Q2:} State the mood of your three claims (A, E, I, or O for each). Is the form valid if written as All M are P; All S are M; therefore All S are P?

\answerbox{Part 1 Q2 ANSWER BOX}{2.25in}

\newpage

\section*{Part 2: Middle Term and Figure}

\textbf{Part 2 Q1:} Identify S, P, and M in Brutus's formal map. Then write the figure pattern using S, P, and M.

\answerbox{Part 2 Q1 ANSWER BOX}{2.35in}

\textbf{Part 2 Q2:} Does Brutus's formal shape by itself prove that Caesar was dangerously ambitious? Explain the difference between valid form and proved premise.

\answerbox{Part 2 Q2 ANSWER BOX}{2.45in}

\newpage

\section*{Part 3: Antony's Reply}

\textbf{Part 3 Q1:} List Antony's three pieces of evidence against the claim that Caesar was dangerously ambitious.

\answerbox{Part 3 Q1 ANSWER BOX}{2.25in}

\textbf{Part 3 Q2:} Does Antony give a cleaner formal syllogism than Brutus, or does he mainly attack Brutus's minor premise? Explain.

\answerbox{Part 3 Q2 ANSWER BOX}{2.45in}

\newpage

\section*{Part 4: Formal Judgment}

\textbf{Part 4 Q1:} Write one possible syllogism for Antony's reasoning. Then name its mood if possible and say whether the conclusion follows from your premises.

\answerbox{Part 4 Q1 ANSWER BOX}{2.75in}

\textbf{Part 4 Q2:} Which speaker gives the stronger form? Which gives stronger evidence against the pressure-point premise? Explain the difference.

\answerbox{Part 4 Q2 ANSWER BOX}{2.25in}

\end{document}
